\documentclass[12pt,letterpaper]{article}
\usepackage[T1]{fontenc}
\usepackage{amsmath,amssymb,amsthm,microtype}
\usepackage{newtxtext,newtxmath}
\usepackage[letterpaper,margin=1in]{geometry}
\usepackage[hidelinks]{hyperref}
\usepackage{enumitem}
\title{The November 2026 RH Formalization and Zeta Summation Operators\\
\large Modern companion to the corrected 1981 edition, Draft 1}
\author{David Joyner\\\texttt{wdjoyner@gmail.com}}
\date{November 28, 2026}
\begin{document}
\maketitle
\begin{abstract}
This note was prepared with assistance from ChatGPT (OpenAI), as described
in the acknowledgments. We record the formal verification of a recently
announced RH proof and identify implications and non-implications for
the zeta-associated summation-operator models. The mathematical narrative
of the corrected 1981 edition does not depend on this modern development.
\end{abstract}
\section{Verification status}
On November 20, 2026 Levent Alp\"oge announced a proof of the
Riemann hypothesis (RH) developed with Claude. The Lean FRO and
Mathlib Initiative certified on November 21 that the formal proof
establishes Mathlib's vetted RH statement and passes independent
kernel and statement-comparison checks \cite{LF26}.
Rubinstein reports independent checks of the original and shortened
formalizations \cite{Ru26}. This is substantial formal verification,
but human-readable mathematical expositions and ordinary review are
still developing. We have not independently repeated these checks.

\section{The principal conditional deductions}
Assuming RH, every nontrivial zero of the ordinary zeta function
has the form
\begin{equation}\label{eq:rho}
 \rho=\tfrac12+i\gamma,\qquad\gamma\in\mathbb R.
\end{equation}
Consequently any upper envelope $\theta$ for real parts of nontrivial
zeta zeros equals $\tfrac12$. In particular, the classical consequences
\begin{equation}\label{eq:prime-error}
 \psi(x)=x+O(x^{1/2}\log^2x),\qquad
 \pi(x)=\operatorname{Li}(x)+O(x^{1/2}\log x)
\end{equation}
become available for ordinary primes. These estimates use the
classical explicit-formula machinery in addition to zero location.

\subsection{The Sobolev-space boundary}
For $S_k$, the powers $y^s$ belong to the algebraic operator
domain precisely when $\Re s>1/(2k)$. Thus for $k\ge2$,
the exponents in \eqref{eq:rho} are admissible, while for $k=1$
none is: the critical line equals the strict boundary.
This is not a technicality. Any assertion that RH yields
zero eigenvectors in $S_1$ overlooks the integrability condition.

\subsection{Joint kernels and generalized eigenvectors}
For the pair of symbols
\begin{equation}\label{eq:pair}
m_q(s)=(1-q^{1-s})\zeta(s),\qquad q=2,3,
\end{equation}
their common zeros have the form \eqref{eq:rho} under RH.
A zero of order $m$ can produce the logarithmic powers
$y^\rho(\log y)^j$, $0\le j<m$, whenever domain membership
is verified. RH does not establish zero simplicity or
spectral synthesis. Equality of algebraic zero-eigenspaces
is weaker than equality of the closures of the operators'
joint kernels.

\subsection{The shifted dilation generator}
The generator $\Delta=y\,d/dy$ satisfies
$\Delta y^\rho=\rho y^\rho$ on its natural algebraic domain.
Under RH, $-i(\Delta-\tfrac12)$ has real eigenvalues $\gamma$
on the admissible zero powers. This does not establish symmetry,
self-adjointness, orthogonality, a Riesz basis, or a Hilbert--P\'olya
operator. Multiple zeros would still create Jordan chains.

\section{Grosswald--Schnitzer perturbations}
For perturbed primes $p_n\le a_n\le p_{n+1}$,
Grosswald--Schnitzer prove that
\begin{equation}\label{eq:quotient}
 \zeta_A(s)=Q_A(s)\zeta(s),\qquad Q_A(s)\ne0
 \quad(\Re s>0)
\end{equation}
with $Q_A$ analytic there \cite{GS78}. Therefore the
nontrivial zeros of $\zeta_A$ in this region have
the same ordinates and multiplicities as those of $\zeta$,
and under RH all lie on $\Re s=\tfrac12$.
The $S_1$ boundary issue remains unchanged.
RH does not imply that $Q_A$ or $Q_A^{-1}$ is bounded
on $\Re s>1/(2k)$, nor that full summation operators
are similar or have identical nonzero spectra.

For endpoint shifts $a_n=p_{n+e_n}$, bounded comparison
still depends on the weighted incidence condition
\begin{equation}\label{eq:boundary}
 \sum_{j\ge1}p_j^{-1/(2k)}
 |(e_j-f_j)-(e_{j-1}-f_{j-1})|<\infty.
\end{equation}
Its validity is arithmetic and independent of RH.
Even infinitesimal prime shifts can produce badly
conditioned multipliers, whereas telescoping endpoint
shifts may remain bounded.

\section{Limits of transfer and research opportunities}
The generalized zeta functions generated by
$\operatorname{Li}^{-1}(n)$ and random Beurling systems
need not possess a quotient $Q_A$ as in
\eqref{eq:quotient}. Their zero sets, branch points,
and possible natural boundaries cannot be inferred
from RH for the ordinary zeta function.
The separately announced Dirichlet-$L$ generalization
does not cover all Beurling, Hecke, Dedekind, or
automorphic zeta functions without further arguments.

Promising research problems include determining the closed joint
kernel beyond algebraic powers, studying a natural replacement
for $S_1$ with critical-line boundary distributions,
classifying full graph domains and pseudospectral stability
of prime-perturbed summation operators, and investigating
whether the formal RH argument supplies any natural operator
structure rather than simply constraining a symbol's zeros.
These are open questions, not automatic consequences of RH.

\section{Glossary}
\begin{description}[style=nextline,leftmargin=1.5em]
\item[Critical line] The line $\Re s=\tfrac12$. RH locates
all nontrivial ordinary zeta zeros on this line.
\item[Formal verification] Proof checking of a precisely
formulated statement by a proof-assistant kernel. It is
different from an accessible human mathematical exposition.
\item[Graph boundary] The signed incidence divergence of
an edge flow. Weighted summability controls bounded comparison
for endpoint prime perturbations.
\item[Joint kernel] The intersection of the null spaces of
several operators. Its analytic closure is not determined merely
by the algebraic zero set.
\item[Operator domain] The set on which a linear operator acts.
Unbounded operator comparisons require domain control in
addition to formal identities.
\item[Zero multiplicity] The vanishing order of a holomorphic
function at a zero. It determines the number of logarithmic
generalized eigenvectors in the formal spectral model.
\end{description}
\section*{Acknowledgments}
ChatGPT (OpenAI) drafted the modern analysis and explanations.
The cited formal RH checks were done by other parties; this
note does not independently verify them. AI revision is not
an independent certification of any novel operator-theoretic
statement. The project revision log records the use of AI.
\begin{thebibliography}{9}
\bibitem{GS78} E. Grosswald and F. J. Schnitzer,
A class of modified $\zeta$ and $L$-functions,
\emph{Pacific J. Math.} \textbf{74} (1978), 357--364.
\bibitem{LF26} Lean FRO and Mathlib Initiative,
\emph{Lean verification of the Riemann Hypothesis},
joint statement, November 21, 2026,
\url{https://leanprover-community.github.io/blog/posts/riemann-hypothesis/}.
\bibitem{Ru26} M. Rubinstein,
\emph{Alp\"oge/Claude proof of the Generalized Riemann Hypothesis:
exposition and verification}, November 2026,
\url{https://www.math.uwaterloo.ca/~mrubinst/alpoge_claude_grh/}.
\end{thebibliography}
\end{document}